% ============================================
% DOC-18 (ZH)
% 中文翻译 — SPEC-935-R356
% ============================================
\section*{\textbf{DOC-18} --- 圈栏官方报告}
\markboth{DOC-18 --- 圈栏官方报告}{}

\noindent \textit{1916 —— 关于阿拉加迪索（Alagadiço）收容营的重构}

\noindent \textbf{虚构重构 —— 未在国家档案馆发现对应原件}
\textbf{1916 —— 叙事文件，并非转录}

\vspace{0.5em}\noindent\rule{\textwidth}{0.5pt}\vspace{0.5em}

\noindent \textit{行政文件的虚构重构}
\textit{公共救济}
\textit{塞阿拉收容营叙事报告}
\textit{1915—1916年度}

\vspace{0.5em}\noindent\rule{\textwidth}{0.5pt}\vspace{0.5em}

\noindent \textbf{营地:} 福塔莱萨附近的阿拉加迪索（Alagadiço）
\textbf{官方容量:} 本重构不提供
\textbf{期间收容人数:} 无精确数字；现有资料记载有数千人被关押
\textbf{死亡人数:} 本叙事不复述无法核实的总数

\noindent \textbf{虚构报告人的声音:}

\noindent \textit{“集中营，官方称为‘政府牲畜圈’，是作为遏制逃荒者流向首府的紧急措施而建立的。实际上，它们成了处于饥饿衰竭状态的人类的仓库。}

\noindent \textit{食物和卫生条件不足以保护被关押者的生命。本重构不为这段措辞赋予死亡率百分比或档案作者。}

\noindent \textit{建议关闭各营地，并采取有效的救济措施。”}

\vspace{0.5em}\noindent\rule{\textwidth}{0.5pt}\vspace{0.5em}

\noindent \textbf{重构编辑说明:}

\noindent 阿拉加迪索（Alagadiço）收容营在1915年大旱期间位于福塔莱萨附近，这一点有历史记录。上述数字、报告人口吻和行政形式均为虚构手法。本段不声称该营地持续开放至1917年；若阿金在该年被转送属于人物的虚构时间线。

\vspace{0.5em}
\noindent\textit{\textbf{↪ 路线:} \rota{DOC-19} \rotasep \rota{MEM-06}}
